\chapter{Moving and periodic sources}\label{ch:moving}
For $Q(x,t)=q(x-Vt)$, direct change of variables gives
\begin{equation}\widetilde Q(k,\omega)=2\pi\widetilde q(k)\delta(\omega-kV).
\label{eq:moving-spectrum}\end{equation}
Multiplication by modulation $g(t)$ broadens the constraint to
$\widetilde Q=\widetilde q(k)\widetilde g(\omega-kV)$. Strong response requires
$\omega_n(k)\simeq kV$, i.e. source speed equals phase velocity, consistent with
moving-CW experiments \cite{li2018}; group velocity instead determines energy
transport after excitation.

For $N$ elements at $x_m=m\Lambda$ with phase delay $m\varphi$,
\begin{equation}A_N(k)=\sum_{m=0}^{N-1}\ee^{-\ii m(k\Lambda-\varphi)}
=\ee^{-\ii(N-1)\theta/2}\frac{\sin(N\theta/2)}{\sin(\theta/2)},
\quad\theta=k\Lambda-\varphi.\label{eq:array-factor}\end{equation}
The infinite limit is a reciprocal-lattice comb
$k=(\varphi+2\pi\ell)/\Lambda$. Scanned pulses add a temporal delay to each
element and couple $k$ and $\omega$. Spectral intersection is necessary but not
sufficient: the physical amplitude contains the overlap
\cref{eq:physical-overlap}. Causality further means that selectivity develops
only after undesired modal transients and the finite-array turn-on field have
decayed according to the measurement metric \cref{eq:settling}.
